\documentclass[../main]{subfiles}

\begin{document}
\chapter*{Preface to the new Typesetting}
\addcontentsline{toc}{chapter}{Editor's Preface}

The book was TeX'd up by the Texromancers, a latexing group. The credits for the typesetting of this book go to:  Aareyan Manzoor, Jonas Hardt, Evelyn Koo, Yohan Wittgenstein, Grisha Taroyan and others.

Here is a link to a dyslexic friendly version of the book:
\url{https://aareyanmanzoor.github.io/assets/books/homotopical-algebra-dyslexic.pdf} . Here is a link to our \LaTeX : \url{https://github.com/AareyanManzoor/Quillen-Homotopical-Algebra}

We added citations and references with hyperlinks. References to e.g. theorems/lemmas in the book are in blue, while citations to the bibliography is
in red. The bibliography also has URLs now, for easy access. Some of the books
in the bibliography had newer editions, so we went with those.

We also added a small index. This is minimalist as of now, if the reader feels more words should be added to it, please let us know. Also please report typos or anything else to us, readers taking the time to help the proofreading is appreciated. This can be reported at our website: \url{https://aareyanmanzoor.github.io/Texromancers.html}.

We changed some notations, particularly for the name of categories. (sets) is now $\mathbf{Set}$ for example. Underlining is also a relic of the typewritter era, so those were replaced with bolded text, same for category names. For example category $\xunderline{C}$ in the original book is now $\mathbf{C}$.

All diagrams are redrawn in tikz or tikzcd.

\section*{Corrections to the text}
The transcription has been checked against the Springer first edition (1967,
Lecture Notes in Mathematics 43). Missing hypotheses, formulas, diagram labels,
and references have been restored; in particular M0--M5 are the model-category
axioms of Chapter I, \S1, while M6 is the additional closed-model axiom of
Chapter I, \S5. The notation changes described above are retained.

Some errors already occur in the first edition. The following are editorial
corrections, rather than restorations of the original printing (page numbers
here are chapter-local).
\begin{itemize}
  \item In I.1.6, the cylinder factorization uses a cofibration followed by a
  trivial fibration, not two trivial maps. In I.1.10 the target of composition
  is $\pi^r(A,C)$; in I.1.14 the fibrant replacement is a trivial cofibration.
  References there and in I.1.12 to the homotopy-comparison lemma now point to
  Lemma~\ref{lem:1.1.5}, not Lemma~\ref{lem:1.1.4}.
  \item In I.2.8 the homotopy has domain $A\times I$. In I.3.8 exactness uses
  $\partial_*(\Omega p)_*=0$, and in I.3.9 the coaction boundary is
  $(0+\id_{\Sigma A})n$. In I.5.4 the lifted homotopy has target $X$,
  and its boundary in the lifting square is $q^{-1}\cdot tpq$.
  The retract diagram in I.5.5 has been given the standard arrow-category directions.
  \item In I.3.16 the Toda formula uses $X$ and $E$, as required by
  $A\to X\to E\to B$; $\alpha^{-1}$ is interpreted through the adjunction
  $[A,\Omega B]\simeq[\Sigma A,B]$.
  \item In II.3.13 the simplicial-group homotopy is defined on
  $G\otimes\Delta(1)$ and its simplex with $i+1$ zeroes has $q-i$ ones.
  The factorization reference is to Lemma~\ref{lem:2.3.3}.
  \item In II.5.2--5.7 cohomological degree $q$ is represented by
  $\Omega^N\Sigma^{q+N}A$, not $\Omega^{q+N}\Sigma^N A$.
  Accordingly $H_M^q(\Sigma X,A)=H_M^{q-1}(X,A)$. This agrees with
  the chain-complex calculation in the proof of Theorem~5:
  $NK(A,q)=A[q]$ (see also the
  \href{https://stacks.math.columbia.edu/tag/0199}{Stacks Project, Lemma 14.23.5}).
  In II.5.10 the sheaf-to-presheaf functor is on abelian sheaves and the
  covariant mapping-space map follows $Z_\bullet\to Y_\bullet$.
  In II.5.11 the cohomology vanishing is for $q>0$, and in II.5.13
  the projective object is a retract of $FSP$, not $FSF$.
  \item In II.6.2 the last internal-Hom target is $Y$.
  In II.6.7 the symmetry of the shuffle product reads
  $\tau(x\wedge y)=(-1)^{pq}(y\wedge x)$.
  In Theorem~6, $X$ is a right module and $Y$ a left module,
  and the graded Tor expression uses $\pi X,\pi Y$ throughout.
  In II.6.14 the factor $i$ is a trivial cofibration.
  The degree-zero sphere module is explicitly $RS^0=R$; in II.6.12
  suspension of a graded $S$-module is again an $S$-module.
  In II.6.17 the group-homology tensor product is derived.
\end{itemize}

\section*{Compiling this edition}
The default accessible edition requires LuaLaTeX, the genuine
\href{https://opendyslexic.org/}{OpenDyslexic} font family (including bold and
italic weights), and the XITS Math font supplied by TeX Live. Install the
OpenDyslexic fonts under their family name \texttt{OpenDyslexic}; no substitute
font is selected. Bold mathematical symbols use XITS Math rather than the text
font, and the page geometry reserves space for two-line running heads.
A complete build also requires Biber compatible with the
installed \texttt{biblatex}, and MakeIndex. From the repository directory run
\begin{quote}
\texttt{latexmk -lualatex -outdir=build main.tex}
\end{quote}
with these programs on \texttt{PATH}; the final PDF is \texttt{build/main.pdf}.
\end{document}